\section{Vector Spaces}

\subsection{\texorpdfstring{$\rbb^n$ and $\cbb^n$}{Rn and Cn}}

\begin{exercise}{1a1}
    Show that $\alpha + \beta = \beta + \alpha$ for all $\alpha, \beta \in \cbb$.
\end{exercise}

\begin{sol}
    Let $\alpha = a + bi$ and $\beta = c + di$.
    Then
    \begin{align*}
        \alpha + \beta & = (a + bi) + (c + di) \\
        & = (a + c) + (b + d)i \\
        & = (c + a) + (d + b)i \\
        & = (c + di) + (a + bi) = \beta + \alpha. \qedhere
    \end{align*}
\end{sol}

\begin{exercise}{1a2}
    Show that $(\alpha + \beta) + \gamma = \alpha + (\beta + \gamma)$ for all $\alpha, \beta, \gamma \in \cbb$.
\end{exercise}

\begin{sol}
    Let $\alpha = a + bi$, $\beta = c + di$, and $\gamma = e + fi$.
	Then
    \begin{align*}
        (\alpha + \beta) + \gamma & = ((a + bi) + (c + di)) + (e + fi) \\
        & = ((a + c) + (b + d)i) + (e + fi) \\
        & = (a + c + e) + (b + d + f)i \\
        & = (a + (c + e)) + (b + (d + f))i \\
        & = (a + bi) + ((c + di) + (e + fi)) = \alpha + (\beta + \gamma). \qedhere
    \end{align*}
\end{sol}

\begin{exercise}{1a3}
    Show that $(\alpha\beta)\gamma = \alpha(\beta\gamma)$ for all $\alpha, \beta, \gamma \in \cbb$.
\end{exercise}

\begin{sol}
    Let $\alpha = a + bi$, $\beta = c + di$, and $\gamma = e + fi$.
	Then
    \begin{align*}
        (\alpha\beta)\gamma & = ((a + bi)(c + di))(e + fi) \\
        & = ((ac - bd) + (ad + bc)i)(e + fi) \\
        & = (ace - bde - adf - bcf) + (acf - bdf + ade + bce)i \\
        & = (ace - adf - bde - bcf) + (acf + ade + bce - bdf)i \\
        & = (a + bi)((ce - df) + (cf + de)i) \\
        & = (a + bi)((c + di)(e + fi)) = \alpha(\beta\gamma). \qedhere
    \end{align*}
\end{sol}

\begin{exercise}{1a4}
    Show that $\gamma(\alpha + \beta) = \gamma\alpha + \gamma\beta$ for all $\alpha, \beta, \gamma \in \cbb$.
\end{exercise}

\begin{sol}
    Let $\alpha = a + bi$, $\beta = c + di$, and $\gamma = e + fi$.
	Then
    \begin{align*}
        \gamma(\alpha + \beta) & = (e + fi)((a + bi) + (c + di)) \\
        & = (e + fi)((a + c) + (b + d)i) \\
        & = (e(a + c) - f(b + d)) + (e(b + d) + f(a + c))i \\
        & = (ea - fb + ec - fd) + (eb + fa + ed + fc)i \\
        & = (ea - fb) + (eb + fa)i + (ec - fd) + (ed + fc)i \\
        & = (e + fi)(a + bi) + (e + fi)(c + di) = \gamma\alpha + \gamma\beta. \qedhere
    \end{align*}
\end{sol}

\begin{exercise}{1a5}
    Show that for every $\alpha \in \cbb$, there exists a unique $\beta \in \cbb$ such that $\alpha + \beta = 0$.
\end{exercise}

\begin{sol}
    Let $\alpha = a + bi$.
	We want to find $\beta = c + di$ such that $\alpha + \beta = 0$.
	This gives
    \[
    	(a + bi) + (c + di) = (a + c) + (b + d)i = 0 + 0i.
    \]
    Solving for $c$ and $d$, we get $c = -a$ and $d = -b$.
	Hence, $\beta = c + di = -a - bi = -\alpha$.

    To see that $\beta$ is unique, suppose there is another $\gamma \in \cbb$ such that $\alpha + \gamma = 0$.
	Then
    \[
    	\beta = \beta + 0 = \beta + (\alpha + \gamma) = (\beta + \alpha) + \gamma = 0 + \gamma = \gamma.
    \]
    Hence, $\beta$ is unique.
\end{sol}

\begin{exercise}{1a6}
    Show that for every $\alpha \in \cbb$ with $\alpha \neq 0$, there exists a unique $\beta \in \cbb$ such that $\alpha\beta = 1$.
\end{exercise}

\begin{sol}
    Let $\alpha = a + bi$ such that $\alpha \neq 0$.
	Then at least one of $a$ or $b$ is nonzero.
	We want to find $\beta = c + di$ such that $\alpha\beta = 1$.
	Then
    \[
    	(a + bi)(c + di) = (ac - bd) + (ad + bc)i = 1 + 0i
    \]
    which gives the system of equations
    \[
        \begin{cases}
            ac - bd = 1 \\
            ad + bc = 0
        \end{cases} \implies
        \begin{cases}
            a^2c - abd = a \\
            b^2c + abd = 0.
        \end{cases}
    \]
    Adding the two equations, we get
    \[
    	(a^2 + b^2)c = a \implies c = \frac{a}{a^2 + b^2} \implies abd + b^2\left(\frac{a}{a^2 + b^2}\right) = 0 \implies d = \frac{-b}{a^2 + b^2}.
    \]
    Hence,
    \[
    	\beta = c + di = \frac{a}{a^2 + b^2} + \frac{-b}{a^2 + b^2}i.
    \]

    If $\gamma \in \cbb$ is another complex number such that $\alpha\gamma = 1$, then
    \[
    	\gamma = \gamma 1 = \gamma (\alpha\beta) = (\gamma\alpha) \beta = 1\beta = \beta.
    \]
    Hence, $\beta$ is unique.
\end{sol}

\begin{exercise}{1a7}
    Show that
    \[
    	\frac{-1 + \sqrt3i}{2}
    \]
    is a cube root of 1.
\end{exercise}

\begin{sol}
    We have
    \begin{align*}
        (-1 + \sqrt 3i)^3 &= (-1)^3 + 3(-1)^2(\sqrt 3i) + 3(-1)(\sqrt 3i)^2 + (\sqrt 3i)^3 \\
        &= -1 + 3\sqrt 3i - 9i^2 - 3\sqrt 3i^3 \\
        &= -1 + 3\sqrt 3i + 9 - 3\sqrt 3i = 8.
    \end{align*}
    Hence,
    \[
        \left(\frac{-1 + \sqrt3i}{2}\right)^3 = \frac{(-1 + \sqrt 3i)^3}{2^3} = \frac{8}{8} = 1. \qedhere
    \]
\end{sol}

\begin{exercise}{1a8}
    Find two distinct square roots of $i$.
\end{exercise}

\begin{sol}
    Let $z = a + bi \in \cbb$ such that $z^2 = i$.
    Then
    \[
        (a + bi)^2 = a^2 - b^2 + 2abi = 0 + 1i,
    \]
    giving the system of equations
    \[
        \begin{cases}
            a^2 - b^2 = 0 \\
            2ab = 1.
        \end{cases}
    \]
    The first equation yields $a = \pm b$.
    If $a = b$, then
    \[
        2b^2 = 1 \implies b^2 = \frac{1}{2} \implies b = \pm\frac{1}{\sqrt 2} \implies a = \pm\frac{1}{\sqrt 2}.
    \]
    If $a = -b$, then $-2b^2 = 1$, which has no real solutions.
    Hence, the two distinct square roots of $i$ are
    \[
    	\frac{1}{\sqrt 2} + \frac{1}{\sqrt 2}i \quad \text{and} \quad -\frac{1}{\sqrt 2} - \frac{1}{\sqrt 2}i. \qedhere
    \]
\end{sol}

\begin{exercise}{1a9}
    Find $x \in \rbb^4$ such that
    \[
    	(4, -3, 1, 7) + 2x = (5, 9, -6, 8).
    \]
\end{exercise}

\begin{sol}
    Let $x = \tuple* x \in \rbb^4$.
	The equation becomes
    \[
    	(4, -3, 1, 7) + 2(x_1, x_2, x_3, x_4) = (5, 9, -6, 8),
    \]
    which simplifies to the system of equations
    \[
        \begin{cases}
            5 = 4 + 2x_1 \\
            9 = -3 + 2x_2 \\
            -6 = 1 + 2x_3 \\
            8 = 7 + 2x_4
        \end{cases} \implies 
        \begin{cases}
            2x_1 = 1 \\
            2x_2 = 12 \\
            2x_3 = -7 \\
            2x_4 = 1 
        \end{cases} \implies
        x = \left(\frac{1}{2}, 6, -\frac{7}{2}, \frac{1}{2}\right). \qedhere
    \]
\end{sol}

\begin{exercise}{1a10}
    Explain why there does not exist $\lambda \in \cbb$ such that
    \[
    	\lambda(2 - 3i, 5 + 4i, -6 + 7i) = (12 - 5i, 7 + 22i, -32 - 9i).
    \]
\end{exercise}

\begin{sol}
    If such a $\lambda$ exists, then we have the system of equations
    \[
        \begin{cases}
            \lambda(2 - 3i) = 12 - 5i \\
            \lambda(5 + 4i) = 7 + 22i \\
            \lambda(-6 + 7i) = -32 - 9i
        \end{cases}
    \]
    The first equation gives
    \[
    	\lambda = \frac{12 - 5i}{2 - 3i} = \frac{(12 - 5i)(2 + 3i)}{(2 - 3i)(2 + 3i)} = \frac{24 + 36i - 10i - 15i^2}{4 + 9} = \frac{24 + 26i + 15}{13} = 3 + 2i.
    \]
    However, using this value for $\lambda$ in the third equation, we get
    \[
    	(3 + 2i)(-6 + 7i) = -18 + 21i - 12i + 14i^2 = -18 + 9i - 14 = -32 + 9i \neq -32 - 9i. \qedhere
    \]
\end{sol}

\begin{exercise}{1a11}
    Show that $(x + y) + z = x + (y + z)$ for all $x, y, z \in \fbb^n$.
\end{exercise}

\begin{sol}
    Let $x = \tuple* x, y = \tuple* y, z = \tuple* z \in \fbb^n$.
    Then
    \begin{align*}
        (x + y) + z &= (\tuple x + \tuple y) + \tuple z \\
        &= \tuple{x, + y} + \tuple z \\
        &= \tuple{(x, + y,) +z} \\
        &= (x_i + (y_i + z_i))_1^n \\
        &= \tuple x + \tuple{y, + z} \\
        &= \tuple x + (\tuple y + \tuple z) = x + (y + z). \qedhere
    \end{align*}
\end{sol}

\begin{exercise}{1a12}
    Show that $(ab)x = a(bx)$ for all $a, b \in \fbb$ and $x \in \fbb^n$.
\end{exercise}

\begin{sol}
    Let $a, b \in \fbb$ and $x = \tuple* x \in \fbb^n$.
    Then
    \[
        (ab)x = (ab)\tuple x = \tuple{(ab)x} = \tuple{a(bx)} = a\tuple{bx} = a(b\tuple x) = a(bx). \qedhere
    \]
\end{sol}

\begin{exercise}{1a13}
    Show that $1x = x$ for all $x \in \fbb^n$.
\end{exercise}

\begin{sol}
    Let $x = \tuple* x \in \fbb^n$.
    Then
    \[
        1x = 1\tuple x = \tuple{1x} = \tuple x = x. \qedhere
    \]
\end{sol}

\begin{exercise}{1a14}
    Show that $\lambda(x + y) = \lambda x + \lambda y$ for all $\lambda \in \fbb$ and $x, y \in \fbb^n$.
\end{exercise}

\begin{sol}
    Let $\lambda \in \fbb$ and $x = \tuple* x, y = \tuple* y \in \fbb^n$.
    Then
    \begin{align*}
        \lambda(x + y) &= \lambda(\tuple x + \tuple y) \\
        &= \lambda\tuple{x, + y} \\
        &= \tuple{\lambda(x, + y)} \\
        &= (\lambda x_i + \lambda y_i)_1^n \\
        &= \tuple{\lambda x} + \tuple{\lambda y} \\
        &= \lambda\tuple x + \lambda\tuple y = \lambda x + \lambda y. \qedhere
    \end{align*}
\end{sol}

\begin{exercise}{1a15}
    Show that $(a + b)x = ax + bx$ for all $a, b \in \fbb$ and $x \in \fbb^n$.
\end{exercise}

\begin{sol}
    Let $a, b \in \fbb$ and $x = \tuple* x \in \fbb^n$.
    Then
    \[
        (a + b)x = (a + b)\tuple x = \tuple{(a + b)x} = \tuple{ax + bx} = \tuple{ax} + \tuple{bx} = a\tuple x + b\tuple x = ax + bx. \qedhere
    \]
\end{sol}